\section{梯度下降法}
牛顿迭代法的核心思想是利用函数的导数信息来构造一个迭代式，使得数列快速收敛到零点。这个思想在优化算法中也有重要应用，特别是在梯度下降法中。梯度下降法的目标是最小化一个函数 $f(x)$，其迭代式为
\[x_{n+1} = x_n - \alpha \nabla f(x_n),
\]其中 $\alpha$ 是学习率，$\nabla f(x_n)$ 是函数 $f$ 在 $x_n$ 处的梯度。可以看出，梯度下降法是牛顿迭代法的一种简化版本，后者在更新时还考虑了    二阶导数信息（即 Hessian 矩阵）。因此，牛顿迭代法在某些情况下收敛速度更快，但计算成本也更高；而梯度下降法则更适合大规模优化问题。   

